% Same code as the original growth problem; qualitative solutions for Fall 2026.
\question Classify the work as linear, sublinear, or superlinear relative to
\lstinline{n}, and explain your reasoning in words. Assume \lstinline{n} is
a non-negative integer and arithmetic takes a fixed amount of work.
If different inputs follow different execution paths, describe them separately.
\begin{parts}

\part
\begin{lstlisting}
def strange_add(n):
    if n == 0:
        return 1
    else:
        return strange_add(n - 1) + strange_add(n - 1)
\end{lstlisting}
\begin{solution}[0.25in]
Superlinear. Each non-base call makes two calls on the same smaller input.
Increasing \lstinline{n} by one adds another branching level and roughly
doubles the total number of calls. The same smaller inputs are computed
repeatedly.

If every recursive call checks a shared cache, each distinct input from 0
through \lstinline{n} needs to be computed only once. A first call with an
empty cache then does linear work. Repeating it with the result already
cached takes a fixed amount of work, which is sublinear. Storing only the
top-level result would not remove repeated work within the first call.
\end{solution}

\part
\begin{lstlisting}
def stranger_add(n):
    if n < 3:
        return n
    elif n % 3 ==  0:
        return stranger_add(n - 1) + stranger_add(n - 2) + stranger_add(n - 3)
    else:
        return n
\end{lstlisting}
\begin{solution}[0.25in]
Linear for positive multiples of 3; sublinear (a fixed amount of work)
otherwise, including \lstinline{n = 0}.

For \lstinline{stranger_add(9)}, the calls on 8 and 7 return immediately;
only the call on 6 continues. At the next stage, the calls on 5 and 4
return immediately; only 3 continues. The continuing chain is
$9 \to 6 \to 3 \to 0$.

Each stage does a fixed amount of work and reduces the continuing input by
3. Doubling a positive multiple-of-3 input roughly doubles the number of
stages. For inputs not divisible by 3, the initial call returns immediately.

Unlike \lstinline{strange_add}, this function does not repeatedly solve the
same inputs during a single call, so caching does not improve its first-call
growth category. A persistent cache can still help repeated top-level calls.
\end{solution}
\end{parts}
